\chapter{Пластичность и ореол}\label{ch:plasticity}

\section{Почему без пластичности нельзя}

\fig[0.85]{F09_Yield}{Несоответствие гидрида и деформация, при которой металл начинает течь.}

Предел упругости циркония при температуре выпадения --- около $\sigma_y/E=350/90\,000\approx0.4\%$
(Zircaloy-4 при 350\,$^\circ$C; у рекристаллизованного Э635 ещё меньше: 226~МПа при 380\,$^\circ$C).
Несоответствие гидрида в 12--18 раз больше (рис.~\ref{fig:F09_Yield}). Упругая теория предсказывает
у кончика пластинки напряжения в гигапаскали; на деле металл вокруг течёт, и поле там совсем
другое. Это видно и по старой модели: выгоду от соседей приходилось обрезать «потолком»
$\pm\sigma_{\text{cap}}$ (90--240~МПа), который подбирался по опыту. Подобранный потолок ---
слабое место, и эта глава о том, чем его заменить.

\section{Пластичность в двух формулах}

\paragraph{Критерий текучести Мизеса.} Металл течёт, когда эквивалентное напряжение
\begin{equation}
\sigma_{\text{экв}}=\sqrt{\tfrac32\,s:s},\qquad s=\sigma-\tfrac13(\mathrm{tr}\,\sigma)\,I
\end{equation}
достигает предела $\sigma_y+h\,p$, где $p$ --- накопленная пластическая деформация, $h$ ---
упрочнение (в расчётах $\sigma_y=350$~МПа, $h=200$~МПа). Давление ($\mathrm{tr}\,\sigma$) на
течение не влияет --- важное свойство, к которому вернёмся в гл.~\ref{ch:limits}.

\paragraph{Пластическая деформация.} Полная деформация делится на упругую, собственную и
пластическую:
\begin{equation}
\eps=\eps^{\text{упр}}+\eps^*+\eps_p,\qquad \sigma=C:\left(\eps-\eps^*-\eps_p\right).
\end{equation}
$\eps_p$ растёт в направлении $s$ (ассоциированное течение) и не меняет объёма
($\mathrm{tr}\,\eps_p=0$). На каждом шаге решения пробное упругое напряжение, вышедшее за
поверхность текучести, «возвращают» на неё --- это алгоритм возврата.

\begin{exercise}
При одноосном растяжении $\sigma=\mathrm{diag}(\sigma_1,0,0)$ покажите, что
$\sigma_{\text{экв}}=|\sigma_1|$. Чему равно $\sigma_{\text{экв}}$ при гидростатическом сжатии
$\sigma=-pI$? Почему металл перед кончиком гидрида может выдерживать гигапаскаль сжатия
без течения?
\end{exercise}

\section{Решатель упругопластичности на Фурье}

Чтобы считать поля с пластичностью быстро и на тех же сетках, что автомат, в проекте написан
свой решатель (\code{fe/hill\_fft.py}). Он устроен так.
\begin{itemize}
\item Периодическая ячейка, неизвестные --- смещения на сетке; производная --- разность вперёд.
\item Нелинейная задача решается методом Ньютона; на каждой итерации линейная система --- методом
  сопряжённых градиентов с предобуславливателем, который точно обращает однородную упругость
  через Фурье (тот же $K^{-1}$ из гл.~\ref{ch:misfit}).
\item Среднее напряжение ячейки задаётся (нагрузка), средняя деформация ищется.
\item Обобщённая плоская деформация: всё однородно по оси трубы, $\eps_{zz}$ --- одно число.
\item Порядок, как в опыте: сначала нагрузка, затем несоответствие пластинки нарастает от 0 до 1.
\end{itemize}
Мера --- \emph{работа превращения} на единицу объёма пластинки:
\begin{equation}
W=\int_0^1\langle\sigma\rangle_{\text{пл}}:\eps^*\,dt,\qquad g=\frac{W}{\eps_n}\ \ [\text{МПа}],
\end{equation}
где $\langle\cdot\rangle_{\text{пл}}$ --- среднее по пластинке, $t$ --- доля превращения. В упругости
$W$ сводится к $\sigma_{\text{прил}}:\eps^*$ минус собственная энергия; с пластичностью это честная
работа с учётом течения.

\begin{idea}[Главный результат пластичности]
Пластичность \emph{втрое} усиливает выбор ориентации нагрузкой: $g_{\text{рад}}-g_{\text{окр}}\approx1.0\,\sigma_\theta$
вместо упругих $0.36\,\sigma_\theta$. У окружной пластинки нагрузка тянет вдоль неё, а расширение
вдоль пластинки ($\eps_t$) снимается течением --- работать не на чем. У радиальной нагрузка тянет
по нормали, где несоответствие снимается упруго, раскрытием граней.
\end{idea}

\begin{exercise}
Предельная оценка: пусть течение полностью снимает несоответствие вдоль пластинки ($\eps_t$ в
плоскости $r$--$\theta$ «не работает»), а по нормали несоответствие остаётся. Чему тогда равны
$g_{\text{окр}}$ и $g_{\text{рад}}$ под $\sigma_\theta$? Сравните разницу с 1.0$\,\sigma_\theta$ из
расчёта.
\end{exercise}

\section{Сверка с методом конечных элементов}

Свой решатель нужно проверить независимым кодом. Проверкой служил CalculiX (МКЭ) на той же
задаче: та же периодическая ячейка $30\times30$~мкм, обобщённая плоская деформация (уравнения
связи для периодичности и управляющий узел для средней $\eps_{zz}$), квадратичные элементы с
шагом 0.04~мкм у границ пластинок --- в 2.5 раза мельче клетки Фурье. Пластинки выделяются по
очереди: нагрузка, затем старые пластинки, затем новая.

\fig[0.85]{F12_FEM}{Фурье против МКЭ для стопок пластинок при 110~МПа: разность «радиальная
минус окружная» и коллективная прибавка (на сколько соседи облегчают выпадение новой пластинки).
Сама работа новой пластинки расходится на 1--21~МПа из $\approx2700$.}

Числа совпадают с точностью, достаточной для всех выводов: колода --- до 2~МПа, цепочка --- 21~МПа
(перемычка 0.8~мкм --- всего 8 клеток Фурье). Поля в матрице совпадают визуально; разница до
$\pm100$~МПа держится в первых клетках у границ и вдоль полос сдвига.

\section{Ореол: пластичность в упругом автомате}

Решать упругопластическую задачу на каждом шаге автомата --- слишком дорого (тысячи пластинок,
сотни шагов). Выручает простое наблюдение.

\begin{keyf}[Теорема об ореоле]
Если модули гидрида и матрицы одинаковы, напряжения упругопластического решения \emph{в точности}
равны упругому решению с собственной деформацией $\eps^*+\eps_p$, где $\eps_p$ --- найденная
пластическая деформация.
\end{keyf}
Доказательство в одну строку: уравнения $\sigma=C:(\eps-\eps^*-\eps_p)$, $\nabla\cdot\sigma=0$,
$\eps=\operatorname{sym}\nabla\mathbf u$ при известном $\eps_p$ линейны и совпадают с упругой
задачей для суммарной собственной деформации; решение единственно.

Значит, каждой пластинке можно приписать её пластическую деформацию --- \emph{ореол} --- и
считать поле старым упругим Фурье, без потолка.

\fig{F10_Halo}{Слева --- зона течения вокруг одной пластинки. В середине и справа --- выгода
соседа той же ориентации без пластичности и с ней. Течение «съедает» острые лепестки у кончиков,
отодвигает выгодное место дальше вперёд и делает невыгодным широкий ореол вокруг пластинки.}

\paragraph{Таблица ореолов.} Ореол зависит от длины пластинки, от угла $\alpha$ между окружным
напряжением и её нормалью и от величины нагрузки. Посчитано 85 вариантов на Фурье:
полудлины 0.6--2.5~мкм, $\alpha=0,30,60,90^\circ$, $\sigma_\theta=0,75,125,175,250$~МПа.
В автомате промежуточные значения берутся интерполяцией; длинную пластинку получают
«растягиванием» ближайшей короче (кончики раздвигаются, середина заполняется); ореол поворачивается
в оси пластинки и усредняется по клеткам. Растущей пластинке ореол обновляется по мере роста.

\paragraph{Складываются ли ореолы?} Пластичность нелинейна, и ореол пластинки, выросшей рядом с
другими, не обязан совпадать с ореолом одиночной. Проверка на стопках из трёх пластинок
(рис.~\ref{fig:F11_Superpose}, таблица~\ref{tab:superpose}).

\fig{F11_Superpose}{Выгода $g$ от соседей (без вклада нагрузки) для колоды и цепочки из трёх
радиальных пластинок под 110~МПа. Слева --- честный упругопластический расчёт, в середине ---
упругий расчёт с суммой ореолов одиночной пластинки, справа --- как было в автомате (упругость с
потолком $\pm180$~МПа). Рамки --- места следующей пластинки.}

\begin{table}[htbp]\centering\small
\caption{$g$ соседей на месте следующей пластинки, МПа.}\label{tab:superpose}
\begin{tabular}{lrrr}\toprule
случай & честный расчёт & ореолы (сетка 0.4~мкм) & потолок $\pm180$\\\midrule
колода, без нагрузки & $-210$ & $-117$ & $+27$\\
колода, радиальные, 110~МПа & $-128$ & $-124$ & $+27$\\
колода, окружные, 110~МПа & $-221$ & $-149$ & $+27$\\
цепочка, без нагрузки & $+304$ & $+244$ & $+180$\\
цепочка, радиальные, 110~МПа & $+375$ & $+338$ & $+180$\\
цепочка, окружные, 110~МПа & $+231$ & $+130$ & $+180$\\\midrule
средняя ошибка & --- & $\approx60$ & $\approx170$\\\bottomrule
\end{tabular}
\end{table}

\begin{idea}
Сложение ореолов не точное, но в три раза лучше потолка и, главное, верно по знаку: место над
колодой с пластичностью \emph{невыгодно}, а потолок делал его выгодным. И ещё одно: под нагрузкой
цепочка радиальных пластинок притягивает следующую сильнее, чем без нагрузки, а цепочка окружных ---
слабее. Нагрузка и соседи складываются нелинейно --- этого механизма в упругом автомате не было.
\end{idea}

\begin{exercise}
Почему ореолы не обязаны складываться? Подумайте, что происходит в перемычке между двумя
пластинками, если их пластические зоны перекрываются, и почему ошибка больше для окружных
пластинок под нагрузкой (подсказка: где нагрузка помогает течению, а где мешает).
\end{exercise}

\begin{warnbox}
Таблица ореолов посчитана для одного материала: Мизес, $\sigma_y=350$~МПа (Zircaloy-4 при
$\approx350\,^\circ$C), изотропно, в двумерном сечении. Для Э635 ($\sigma_y\approx226$~МПа) нужна
своя таблица; предел текучести при остывании растёт, а таблица от температуры не зависит.
Кристаллическая пластичность (системы скольжения зерна) в ореол не входит.
\end{warnbox}

\begin{codebox}
\textbf{Где в коде.} \code{fe/hill\_fft.py} --- решатель (Мизес и Хилл); \code{fe/stack\_runs.py},
\code{fe/stack\_fe.py}, \code{fe/stack\_compare.py} --- стопки и сверка с CalculiX;
\code{fe/halo\_runs.py}, \code{fe/halo\_table.py} --- таблица ореолов (\code{data\_halo/halo\_tab.npz});
\code{fe/halo\_check.py} --- проверка сложения; \code{halo.py} и опция \code{KParams.halo} --- ореолы в автомате.
\end{codebox}
